\begin{solution}
\begin{align*}
B\pi r^2 &= n\left(\frac{h}{e}\right) \\
r^2 &= \frac{nh}{\pi eB} \\
\intertext{For circular motion of the electron in the magnetic field,}
e v B &= \frac{mv^2}{r} \\
v &= \frac{eBr}{m} \\
\intertext{Magnetic moment of the revolving electron is}
\mu &= IA = \left(\frac{ev}{2\pi r}\right)\left(\pi r^2\right) \\
      &= \frac{evr}{2} \\
      &= \frac{e}{2}\left(\frac{eBr}{m}\right)r \\
      &= \frac{e^2Br^2}{2m} \\
      &= \frac{e^2B}{2m}\cdot \frac{nh}{\pi eB} \\
      &= \frac{nhe}{2\pi m} \\
\intertext{In the lowest energy state, the smallest allowed integer is $n = 1$. Hence,}
\mu &= \frac{he}{2\pi m} \\
\intertext{Therefore, the correct option is (d).}
\end{align*}
\end{solution}
